\documentclass[11pt]{article}

\usepackage{amsmath,amssymb,graphicx,longtable}
\usepackage[margin=1in]{geometry}
\usepackage{booktabs}
\usepackage{array}

\title{Limits of Frozen Pre-Geometry on a Kernel-Induced Cochain Operator Architecture}
\author{Tomislav Rupi\'c \\ Independent Researcher}
\date{March 2026}

\begin{document}

\maketitle

\begin{abstract}
This paper consolidates the full frozen-architecture arc of the HAOS-IIP program through Stage 18A. The sequence begins with construction and validation of a kernel-induced cochain operator substrate, then moves through packet propagation, morphology atlases, collective interaction structure, coarse-grained envelope organization, scale-coupled fragility, weak binding probes, relational-ordering diagnostics, relational-feedback tests, phase-synchronization tests, and finally a minimal retarded-memory probe. Taken stage by stage, the program shows that the frozen architecture can support structured morphology, localized ordering, and visually sharpened coarse organization. Taken as a whole, it also shows a clear boundary: weak clean extensions do not promote those structures into a robust collective law. No weak effective binding channel appears, no geometry-like back-reaction appears, no weak synchronization channel appears, and no minimal retarded-memory channel appears. The bounded conclusion is therefore simple: frozen pre-geometry can organize, but it does not bind. This is not a null program. It is a boundary program. Its value lies in sharply delimiting what a frozen interaction-invariant substrate can and cannot do, thereby motivating a qualitatively different Phase II regime in which the interaction kernel itself becomes dynamical.
\end{abstract}

\section{Introduction}

The HAOS-IIP sequence through Stage 18A is best read as a single falsification-driven program rather than as a collection of disconnected experiments. The underlying architecture is intentionally conservative. A kernel-induced cochain operator is constructed and then frozen. Subsequent stages do not modify the kernel, projector, or base packet family unless a stage is explicitly defined as an extension test. This discipline makes negative results meaningful. If a new organizing channel does not appear, the failure cannot be attributed to a drifting substrate.

This paper packages that frozen arc as a Phase I result. The aim is not to claim that pre-geometric behavior fails in all possible models. The aim is narrower and stronger: to determine what this frozen kernel-induced cochain architecture can support before any substrate back-reaction is introduced. The resulting map is remarkably coherent. Structured morphology appears. Multi-packet organization appears. Coarse basin structure appears. Localized metric-like ordering can even be read out diagnostically. But the architecture does not naturally convert these forms of organization into an emergent binding law, an effective geometric back-reaction, or a stable collective timing channel under the weak extensions tested.

The thesis line of Phase I is therefore operational rather than metaphysical:
\begin{quote}
Frozen pre-geometry can organize, but it does not bind.
\end{quote}

\section{Scope of the Frozen Program}

For clarity, the notation ``0--18'' in this paper denotes the full Phase I sequence from the operator-architecture foundation through the Stage 18A retarded-memory probe. The later stage numbers (10--18A) correspond to the explicit pre-geometric atlas ladder; the earlier portion of the sequence is the architectural and wave-packet foundation captured in papers 00.1 through 05.1.

The entire Phase I program shares four constraints:
\begin{itemize}
\item the kernel-induced operator architecture is frozen,
\item the projected transverse sector remains the main dynamical sector,
\item packet comparisons are made on the same coherent seed family whenever possible,
\item all extension stages are falsification-first and use null-control baselines.
\end{itemize}

These constraints are what allow the final boundary statement to carry weight.

\section{Stage-by-Stage Arc}

Table~\ref{tab:arc} compresses the complete Phase I sequence into its main technical questions and bounded outcomes.

\begin{longtable}{>{\raggedright\arraybackslash}p{0.12\textwidth} >{\raggedright\arraybackslash}p{0.27\textwidth} >{\raggedright\arraybackslash}p{0.25\textwidth} >{\raggedright\arraybackslash}p{0.24\textwidth}}
\toprule
Sequence & Main question & Positive boundary result & Remaining limit \\
\midrule
\endfirsthead
\toprule
Sequence & Main question & Positive boundary result & Remaining limit \\
\midrule
\endhead
00--04 foundations & Can a kernel-induced cochain operator architecture be constructed, validated, and spectrally organized? & The operator substrate, transverse band, continuum reconstruction diagnostics, Dirac--K\"ahler factorization checks, and gauge-like diagnostics are frozen into a coherent baseline. & This only establishes the substrate; it does not yet establish dynamical collective law. \\
05 wave packets & Can packet propagation remain numerically stable and sector-controlled on the frozen substrate? & Wave-packet propagation and constraint control are established cleanly on the frozen projected sector. & Stability of packet motion does not imply morphology richness or emergent interaction. \\
06 baseline morphology & What single-packet morphology classes exist on the frozen substrate? & A morphology instrument is established, with coherent, dispersive, diffusive, fragmenting, and irregular classes. & This remains single-packet organization only. \\
07 resilience & Are those morphology classes stable under weak perturbation? & Several classes remain identifiable under controlled perturbations. & Resilience does not yet reveal how classes interact or transform collectively. \\
08 transition structure & How do morphology classes deform or flip under controlled transitions? & Transition corridors and class-flip structure become measurable. & Transition structure still does not produce collective law. \\
09 integrated atlas & Can baseline, resilience, and transition structure be packaged as a unified pre-geometric atlas? & Yes. The morphology instrument becomes a reproducible diagnostic layer. & The atlas still concerns isolated packet morphology. \\
10 collective interaction & What appears when coherent packets interact rather than propagate in isolation? & Moderate-separation packets occupy transient-binding corridors; tight-overlap cases can form metastable composite envelopes. & The strongest composites remain morphology, not robust interaction law. \\
11 coarse envelopes & Do collective fields produce stable coarse-grained envelope and basin structure? & Compact clusters can coarse-grain into basin-dominated envelope organization; chain-like cases fluctuate; wider arrangements remain diffuse. & Basin structure appears, but not yet as a stable governing field. \\
12 scale coupling & Do interaction and coarse labels survive coupled refinement? & Coarse basin organization consistently sharpens under refinement. & No scale-robust composite class appears; width coupling remains a major fragility channel. \\
13 weak binding tests & Can minimal weak structural extensions generate an effective binding channel? & The answer is negative across weak envelope feedback, phase-lock retention, and envelope-mediated coupling. & No weak binding mechanism emerges. \\
14 relational ordering & Can frozen collective data be compressed into a metric-like relational ordering? & Localized metric-like ordering appears in some regimes, especially simple corridors or symmetric phase cases. & Ordering is regime-dependent, not composite-anchored, and transport-delay ordering is weak. \\
15 relational feedback & Can derived relational ordering act back on the dynamics? & The answer is negative: ordering is diagnostic bookkeeping, not an effective dynamic organizer. & No geometry-like back-reaction appears. \\
16 phase synchronization & Can weak phase coupling open a collective timing channel? & The answer is negative across all tested representatives. & No self-sustained timing law appears. \\
17 memory / retarded coupling & Can a minimal single-lag retarded term generate organization that instantaneous weak channels could not? & The delayed term produces measurable correlation shifts. & Those shifts remain morphologically inert: no composite gain, no basin gain, no new ordering class. \\
\bottomrule
\caption{Phase I arc through Stage 18A. The value of the sequence lies in the cumulative boundary it imposes on frozen pre-geometry.}
\label{tab:arc}
\end{longtable}

\section{From Morphology to Boundary}

The important structural fact is that the sequence does not merely collect null results. It accumulates a patterned distinction among different layers of organization.

\subsection{Morphology exists}

By the end of the single-packet and collective atlas stages, the frozen architecture clearly supports organized morphology. Packets do not simply dissolve into trivial noise. Single-packet classes are reproducible, collective overlap corridors exist, and tight-overlap cases can generate metastable composite envelopes. This already rules out the trivial hypothesis that the frozen operator substrate is dynamically featureless.

\subsection{Coarse organization exists}

Stage 12 shows that collective organization survives coarse-graining in a real but limited sense. Compact clusters can become basin-dominated; looser or corridor-like cases remain diffuse. This is a stronger statement than visual packet overlap. It means the system can support envelope-level summaries that remain meaningful across the simulation window.

\subsection{Ordering exists diagnostically}

Stage 15 shows that some relational order can be read out from the data. Certain corridor or phase-ordered regimes admit localized metric-like ordering under persistence or phase-correlation diagnostics. This result matters because it shows that geometry-like bookkeeping is not entirely absent. But it remains diagnostic rather than universal, and it is not strongest on the composite anchor.

\subsection{Collective law does not appear}

The critical negative chain begins once weak organizing channels are tested directly. Stage 14 shows that weak structural branches do not create an effective binding channel. Stage 16 shows that derived relational ordering does not feed back into dynamics. Stage 17 shows that weak phase synchronization does not create a timing regime. Stage 18A shows that even a physically motivated minimal retarded term produces only delayed-response correlation shifts, not a dynamical well.

The entire Phase I boundary rests on this transition:
\begin{itemize}
\item organization appears,
\item but organization does not promote itself into law.
\end{itemize}

\section{Why Stage 18A Matters}

Stage 18A is the natural closure point of Phase I because it tests the last clean weak extension before the substrate itself must change. Earlier negative results could still invite a tempting response: perhaps the architecture simply needs memory. A delayed term is the smallest way to test that thought without abandoning the frozen ontology.

The result is clean. Stage 18A adds temporal offset and weak history sensitivity, and the system does register that addition statistically through small delayed-response correlation shifts. But the changes remain below the regime level. No composite lifetime increases. No binding persistence increases. No basin persistence increases. No new ordering class appears. No representative is promoted to refinement.

This is the right kind of negative. It closes a plausible loophole without over-engineering the model. It also clarifies the difference between a delayed ghost and a dynamic well. A fixed lag can replay a weak image of prior packet passage. It cannot, on a frozen substrate, carve the adaptive trench needed to trap later motion.

Within the tested extension class, the frozen linear background therefore acts as a practical ceiling on emergent binding and geometry-like back-reaction.

\section{Why Richer Memory Is Not the Next Step}

After Stage 18A, it is tempting to remain in the same ontological regime and add richer memory structures: state-dependent lag, multiscale history weighting, distributed kernels, or other tuned history mechanisms. This paper rejects that path as the immediate continuation of Phase I.

The reason is methodological. Once increasingly parameterized memory mechanisms are introduced on a frozen background, the question quietly changes. The program ceases to ask whether the architecture naturally supports emergent binding and begins to ask how many nonlinear features can be added before something looks like binding. At that point, apparent success becomes difficult to defend as emergence rather than construction.

The discipline of the HAOS-IIP program has been to keep mechanisms minimal, interpretable, and falsifiable. The sequence through Stage 18A earns its force precisely because it resists parameter wandering. Preserving that integrity requires a sharper transition, not more intricate tuning inside the same frozen ontology.

\section{Phase I Conclusion}

The full frozen arc supports a strong but carefully bounded claim:
\begin{quote}
Frozen pre-geometry can organize, but it does not bind.
\end{quote}

That line condenses the entire Phase I program.

It does not mean the architecture is barren. On the contrary, it means the architecture is rich enough to produce reproducible morphology, collective envelopes, and localized ordering while remaining conservative enough not to confuse those structures with a genuine emergent collective law.

This is why the sequence is best understood as a boundary program. Its value lies in mapping a limit cleanly:
\begin{itemize}
\item structured form can appear,
\item localized relational order can appear,
\item but no weak local, relational, synchronization, or minimal retarded extension generates robust self-binding or geometry-like back-reaction.
\end{itemize}

\section{Phase II Transition: Unfreezing the Kernel}

Once the frozen architecture has been pushed through this boundary, the next move should not be another weak tweak on the same substrate. It should be a regime shift.

Phase II therefore opens with the hypothesis that if emergent relational geometry exists in this framework at all, the interaction substrate must itself become a dynamical degree of freedom. In operational form, the kernel becomes
\[
K(x,y) \rightarrow K(x,y,t).
\]

To preserve falsifiability, that unfreezing must remain tightly bounded:
\begin{itemize}
\item kernel updates must arise from a strict local stress or amplitude functional,
\item substrate relaxation must be slower than packet propagation,
\item the same representatives and early-stop logic must remain in force.
\end{itemize}

This is not a rescue patch. It is a new architectural regime. If a bounded dynamic kernel still fails to produce composite persistence gain or basin stabilization, then the deficit lies deeper than substrate frozenness alone. That would be an equally strong result.

\section{Conclusion}

Stages 0--18 form a complete and defensible Phase I arc. The program begins with operator construction, advances through packet morphology and collective organization, and then systematically tests whether weak clean extensions can convert that organization into robust collective law. They do not.

That is the point. The result is not that nothing happened. The result is that a frozen interaction-invariant substrate has now been mapped to a clear structural boundary. It can support organization, but not binding. The right continuation is therefore not more tuning within the frozen ontology. The right continuation is the onset test for substrate back-reaction.

\end{document}
